\documentclass[12pt]{jlreq}
\usepackage{amsmath, amssymb}

\newcommand{\C}{\mathbb{C}}
\newcommand{\R}{\mathbb{R}}
\newcommand{\Z}{\mathbb{Z}}
\newcommand{\N}{\mathbb{N}}
\newcommand{\F}{\mathbb{F}}
\newcommand{\Prob}{\mathbb{P}}
\newcommand{\Exp}{\mathbb{E}}
\newcommand{\dd}{\,d}
\newcommand{\Ker}{\operatorname{Ker}}
\newcommand{\Impart}{\operatorname{Im}}
\newcommand{\Hom}{\operatorname{Hom}}

\begin{document}

単位円板 \(\Delta = \{ z \in \C \mid |z| < 1 \}\) 上で定義された単葉な正則関数 \(f(z)\) で
\[
  f(0) = 0,\quad f'(0) = 1
\]
を満たすものの全体のなす集合を \(\mathcal{S}\) と書く．次の問いに答えよ．

(1) 複素数 \(a \in \Delta\) に対して
  \[
    g(z) = \frac{z}{(1 - az)^2}
  \]
  は \(\mathcal{S}\) に含まれることを示せ．さらに
  \[
    g(z^2) = (G(z))^2,\quad z \in \Delta
  \]
  を満たす \(G \in \mathcal{S}\) が存在することを示せ．

(2) 任意の \(f \in \mathcal{S}\) に対して
  \[
    f(z^2) = (F(z))^2,\quad z \in \Delta
  \]
  を満たす \(F \in \mathcal{S}\) が存在することを示せ．

\end{document}